\documentclass[11pt,a4paper]{article}
\usepackage[top=1.7cm,bottom=1.5cm,left=1.5cm,right=1.5cm,headsep=8pt,footskip=16pt]{geometry}
\usepackage{enumitem}
\usepackage{amssymb}
\usepackage{array}
\usepackage[table]{xcolor}
\usepackage{fancyhdr}
\usepackage{lastpage}
\usepackage{multicol}
\usepackage{polyglossia}
\setdefaultlanguage[calendar=gregorian, numerals=maghrib]{arabic}
\setotherlanguage{english}
\IfFontExistsTF{Amiri}{%
  \newfontfamily\arabicfont[Script=Arabic, Scale=1.08]{Amiri}
}{%
  \newfontfamily\arabicfont[Script=Arabic, Scale=1.08, Path=../fonts/, Extension=.ttf,
    UprightFont=*-Regular, BoldFont=*-Bold, ItalicFont=*-Italic,
    BoldItalicFont=*-BoldItalic]{Amiri}
}
\IfFontExistsTF{AlBattar}{%
  \newfontfamily\khatt[Script=Arabic, Scale=1.15]{AlBattar}
}{%
  \newfontfamily\khatt[Script=Arabic, Scale=1.15, Path=../fonts/,
    BoldFont=ae_AlBattar.ttf]{ae_AlBattar.ttf}
}
\renewcommand{\headrulewidth}{0.8pt}
\renewcommand{\footrulewidth}{0.8pt}
\fancyhf{}
\rhead{\khatt المدرسة العليا للأساتذة ببوزريعة}
\lhead{\khatt ملحقة بومرداس}
\lfoot{\khatt الأستاذ: د. بن عبد الرزاق يوسف}
\cfoot{\thepage/\pageref{LastPage}}
\pagestyle{fancy}
\parindent 0pt
\setlength{\columnsep}{0.6cm}
\setlength{\columnseprule}{0.4pt}
\setlist[enumerate,1]{label=\textbf{\arabic*.}, leftmargin=*, itemsep=5pt, topsep=2pt}
\begin{document}

\linespread{1.08}\selectfont
\begin{center}
{\LARGE\khatt\textbf{الاختبار السابع والعشرون في الذكاء الاصطناعي}}\\[4pt]
{\normalsize وضعيات إشكالية ــ أسئلة متعددة الاختيارات ــ الفصول من 1 إلى 6}
\end{center}
\vspace{-2pt}\rule{\textwidth}{0.6pt}\\[2pt]
\textbf{التعليمة:} يتضمن الاختبار خمس وضعيات إشكالية. اقرأ كل وضعية بعناية، ثم اختر التحليل أو التصرف الأنسب
من بين الاختيارات الخمسة؛ لكل وضعية إجابة صحيحة واحدة فقط.
اكتب \textbf{حرف} الإجابة الصحيحة في جدول الإجابات في آخر الاختبار.
لكل وضعية 4 نقاط (المجموع 20 نقطة).
\vspace{4pt}

\noindent
{\large\khatt\textbf{الوضعية 1}}\hfill{\small(الفصل 1)}\par\nopagebreak\vspace{3pt}
\setlength{\fboxsep}{7pt}\colorbox{gray!10}{\parbox{\dimexpr\linewidth-14pt}{في لقاء مع الأولياء، سأل أبٌ قلق: «هل هذه الأدوات تفكر مثلنا؟ ابني يقول إنها تعرف كل شيء». وطلب المدير من أستاذة اللغة العربية تقديم تعريف دقيق ومبسّط للذكاء الاصطناعي التوليدي، يطمئن الأولياء دون تهوين، ويوضح فائدته وحدوده. وعدد من الأولياء يشاركونه القلق نفسه. أي تعريف هو الأنسب؟}}\par\nopagebreak
\begin{itemize}[leftmargin=1.0cm, labelsep=6pt, itemsep=3pt, topsep=4pt, parsep=0pt]
  \item[\textbf{أ)}] برامج تعلّمت أنماط اللغة من نصوص كثيرة فتنتج محتوى جديدًا، وهي لا تخطئ إذا كان السؤال واضحًا، فيكفي أن يتعلم الأبناء صياغة الأسئلة الواضحة ليعتمدوا عليها دون مراجعة، وتقل حاجتهم إلى مساعدة الأولياء.
  \item[\textbf{ب)}] محركات بحث متطورة تعرض صفحات جاهزة من الإنترنت كما هي، ولا تنتج أي محتوى جديد، ولذلك فهي لا تخطئ أبدًا، لأن كل ما تعرضه منشور سابقًا في مواقع معروفة، ويتحقق الأولياء من مصادرها بسهولة.
  \item[\textbf{ج)}] عقول إلكترونية تفكر مثل الإنسان وتشعر وتفهم كل ما تقوله، وتعرف كل شيء لأنها قرأت الإنترنت كله، فيمكن للأبناء الاعتماد عليها في كل شيء دون مراجعة من الأولياء أو الأساتذة.
  \item[\textbf{د)}] برامج خطيرة تنسخ واجبات الآخرين وتعطيها للتلاميذ، ولا فائدة تعليمية لها إطلاقًا، والحل الوحيد منع الأبناء من استعمالها منعًا كاملًا في البيت والمدرسة دون أي توجيه.
  \item[\textbf{هـ)}] برامج تعلّمت من كميات كبيرة من النصوص والصور أنماط اللغة والمحتوى، فتنتج محتوى جديدًا عند الطلب؛ تفيد في الشرح والتدريب، لكنها قد تخطئ بثقة، فتحتاج إلى توجيه ومراجعة من الكبار.
\end{itemize}
\par\vspace{6pt}
\noindent
{\large\khatt\textbf{الوضعية 2}}\hfill{\small(الفصل 2)}\par\nopagebreak\vspace{3pt}
\setlength{\fboxsep}{7pt}\colorbox{gray!10}{\parbox{\dimexpr\linewidth-14pt}{أراد فريق بحثي تدريب نموذج يصنّف أبيات الشعر حسب «العاطفة الغالبة»: حزن، فخر، حنين، غضب. كلّف ثلاثة نقاد بتصنيف الأبيات نفسها، فاختلفوا في ربعها تقريبًا، لأن بعض الأبيات تجمع أكثر من عاطفة. ويريد الفريق بيانات صالحة للتدريب. ويخشى الفريق أن يتعلم النموذج من بيانات مضطربة. ما الإجراء الأنسب؟}}\par\nopagebreak
\begin{itemize}[leftmargin=1.0cm, labelsep=6pt, itemsep=3pt, topsep=4pt, parsep=0pt]
  \item[\textbf{أ)}] يقرّ بأن التصنيف تأويلي في جزء منه، فيضع دليلًا بمعايير وأمثلة، ويسمح بتصنيف متعدد للأبيات المركّبة أو يسجّل درجة الاتفاق، ولا يفرض على النموذج يقينًا غير موجود في البيانات.
  \item[\textbf{ب)}] يترك النموذج يصنّف الأبيات بنفسه دون أي تسميات بشرية، لأن النموذج أكثر موضوعية من النقاد، وتصنيفه الآلي للعواطف لا يتأثر بأي تأويل أو اختلاف في وجهات النظر، ويوفر الفريق جهد التصنيف كله.
  \item[\textbf{ج)}] يعتمد تصنيف الناقد الأكبر سنًا وحده، لأن الخبرة الطويلة تلغي الطابع التأويلي للتصنيف، فيصبح لكل بيت عاطفة واحدة صحيحة يتعلمها النموذج دون أي غموض أو اختلاف.
  \item[\textbf{د)}] يختار لكل بيت مختلف فيه عاطفة عشوائية من بين اقتراحات النقاد، لأن العشوائية توزع الأخطاء بالتساوي، فلا تؤثر في تعلّم النموذج ولا في دقته على الأبيات الجديدة.
  \item[\textbf{هـ)}] يحذف كل الأبيات المختلف فيها نهائيًا، ويدرّب النموذج على الأبيات المتفق عليها فقط، ثم يستعمله لتصنيف كل الشعر العربي، بما فيه الأبيات المركّبة، دون أي تنبيه إلى حدوده، فيتعلم النموذج من بيانات نظيفة تمامًا.
\end{itemize}
\par\vspace{6pt}
\noindent
{\large\khatt\textbf{الوضعية 3}}\hfill{\small(الفصل 3)}\par\nopagebreak\vspace{3pt}
\setlength{\fboxsep}{7pt}\colorbox{gray!10}{\parbox{\dimexpr\linewidth-14pt}{طلبت أستاذة من الأداة تبسيط نص علمي عن «التصحّر» لتلاميذ السنة الأولى متوسط، فجاء التبسيط سهلًا، لكنه حذف المصطلحات العلمية المقررة كلها واستبدلها بعبارات عامة، مع أن الدرس يستهدف اكتسابها. وتريد نصًا مبسطًا يحافظ على الهدف. وقد لاحظت أن التلاميذ يحتاجون إلى هذه المصطلحات في الاختبار. أي طلب أنسب؟}}\par\nopagebreak
\begin{itemize}[leftmargin=1.0cm, labelsep=6pt, itemsep=3pt, topsep=4pt, parsep=0pt]
  \item[\textbf{أ)}] «بسّط النص أكثر واحذف كل كلمة صعبة»، لأن التبسيط الكامل هو هدف الدرس، والمصطلحات العلمية يتعلمها التلاميذ لاحقًا في المرحلة الثانوية، فلا حاجة إلى الإبقاء عليها الآن، ويُسعدهم بنص خالٍ من كل صعوبة.
  \item[\textbf{ب)}] «بسّط النص لتلاميذ السنة الأولى متوسط في جمل قصيرة، مع الإبقاء على المصطلحات التالية كما هي: …، وشرح كل مصطلح عند أول وروده بعبارة بسيطة، ثم أضف مسردًا قصيرًا في آخر النص».
  \item[\textbf{ج)}] «بسّط النص واكتب المصطلحات العلمية بالإنجليزية بين قوسين فقط»، لأن المصطلح الإنجليزي أوضح للتلاميذ من المصطلح العربي، ويغني عن شرحه أو إدراجه في مسرد.
  \item[\textbf{د)}] «أعد النص الأصلي كما هو دون تبسيط»، لأن أي تبسيط يحذف المصطلحات بالضرورة، والتلاميذ يتعلمون المصطلحات أفضل عندما يواجهون النص العلمي الأصلي دون مساعدة.
  \item[\textbf{هـ)}] «بسّط النص وأبقِ على المصطلحات»، دون تحديدها، ثم تعتمد النص الجديد دون مراجعة، لأن الأداة تعرف المصطلحات المقررة في كل مستوى دراسي دون أن يذكرها المستخدم، فتوفر وقتًا ثمينًا في التحضير اليومي.
\end{itemize}
\par\vspace{6pt}
\noindent
{\large\khatt\textbf{الوضعية 4}}\hfill{\small(الفصل 5)}\par\nopagebreak\vspace{3pt}
\setlength{\fboxsep}{7pt}\colorbox{gray!10}{\parbox{\dimexpr\linewidth-14pt}{كتبت طالبة في مذكرتها معلومة عن تاريخ تأسيس مجمع اللغة العربية، وأحالت إلى «\textenglish{ChatGPT}» مصدرًا. وطلب منها المشرف تغيير الإحالة. فاقترحت زميلتها الإحالة إلى ويكيبيديا بدلًا منه، لأن «المعلومة نفسها موجودة هناك». وموعد إيداع المذكرة بعد أسبوع، والمشرف صارم في التوثيق، وتريد حلًا نهائيًا لهذه الإحالة. ما التصرف المنهجي الأنسب؟}}\par\nopagebreak
\begin{itemize}[leftmargin=1.0cm, labelsep=6pt, itemsep=3pt, topsep=4pt, parsep=0pt]
  \item[\textbf{أ)}] تحذف الإحالة تمامًا، لأن المعلومات التاريخية العامة لا تحتاج إلى توثيق في المذكرات الجامعية، والمشرف يطلب تغيير الإحالة لا إثباتها، فيكفي حذفها لإرضائه.
  \item[\textbf{ب)}] تُبقي الإحالة إلى الأداة، لأن الأداة جمعت المعلومة من مصادر كثيرة، فهي أشمل من أي مصدر منفرد، والمشرف يعترض على الأدوات الحديثة فقط لأنه لا يعرف طريقة عملها، فيتفهم ذلك في النهاية.
  \item[\textbf{ج)}] تحيل إلى ويكيبيديا كما اقترحت زميلتها، لأن الموسوعة مصدر علمي محكّم في كل الحالات، والإحالة إليها مقبولة في المذكرات الجامعية دون الرجوع إلى المراجع التي تستند إليها، كما يفعل كثير من الطلبة.
  \item[\textbf{د)}] تستعمل الأداة وويكيبيديا للوصول إلى المعلومة ومراجعها، ثم ترجع إلى مصدر موثوق أصيل، كالموقع الرسمي للمجمع أو دراسة محكّمة، وتتحقق منه وتحيل إليه، لأن الأداة والموسوعة المفتوحة نقطتا انطلاق.
  \item[\textbf{هـ)}] تحيل إلى الأداة وويكيبيديا معًا، لأن الجمع بين مصدرين غير محكّمين يعادل مصدرًا محكّمًا واحدًا، فتصبح المعلومة موثقة علميًا دون الرجوع إلى أي مصدر أصيل.
\end{itemize}
\par\vspace{6pt}
\noindent
{\large\khatt\textbf{الوضعية 5}}\hfill{\small(الفصل 6)}\par\nopagebreak\vspace{3pt}
\setlength{\fboxsep}{7pt}\colorbox{gray!10}{\parbox{\dimexpr\linewidth-14pt}{يريد أستاذ تحويل رواية قصيرة مقررة إلى كتاب مسموع بأداة توليد صوتي، بأصوات مختلفة للشخصيات، لمساعدة التلاميذ ضعاف القراءة. جاءت النتيجة جذابة، لكن بعض الكلمات غير المشكولة نُطقت خطأ، وتغيّر صوت إحدى الشخصيات في منتصف الرواية. ويريد توزيع النسخة على ثلاثة أقسام. ما الخطة الأنسب لإنتاج نسخة صالحة؟}}\par\nopagebreak
\begin{itemize}[leftmargin=1.0cm, labelsep=6pt, itemsep=3pt, topsep=4pt, parsep=0pt]
  \item[\textbf{أ)}] يوزّع التسجيل كما هو، لأن الأخطاء القليلة لا تؤثر في الفهم العام، وتغيّر صوت الشخصية يضيف تشويقًا، والتلاميذ ضعاف القراءة لا ينتبهون إلى أخطاء النطق في الكلمات غير المشكولة.
  \item[\textbf{ب)}] يحذف الحوار ويحوّل السرد وحده إلى صوت واحد، لأن تعدد الأصوات هو سبب الأخطاء كلها، والصوت الواحد ينطق الكلمات غير المشكولة نطقًا صحيحًا دائمًا دون حاجة إلى ضبط.
  \item[\textbf{ج)}] يضبط بالشكل الكلمات الملتبسة، ويثبّت إعداد الصوت لكل شخصية في كل المقاطع، ويستمع إلى التسجيل كاملًا لتصحيح المواضع الخاطئة، ويخبر التلاميذ بأن الأصوات مولَّدة.
  \item[\textbf{د)}] يتخلى عن الكتاب المسموع، لأن التلاميذ ضعاف القراءة يجب أن يقرؤوا الرواية بأنفسهم فقط، والاستماع يعطل تعلّمهم القراءة، حتى لو كان التسجيل سليمًا ومصحوبًا بالنص المكتوب.
  \item[\textbf{هـ)}] يضبط الكلمات الملتبسة ثم يوزّع التسجيل دون الاستماع إليه كاملًا، لأن الضبط يصحح النطق ويثبّت الأصوات تلقائيًا، فلا حاجة إلى مراجعة مقاطع الرواية واحدًا واحدًا.
\end{itemize}
\par\vspace{6pt}
\begin{center}
{\khatt\textbf{جدول الإجابات}}\\[3pt]
\renewcommand{\arraystretch}{1.6}
\begin{tabular}{|>{\centering\arraybackslash}m{1.9cm}|*{5}{>{\centering\arraybackslash}m{1.2cm}|}}
\hline
\textbf{رقم السؤال} & \textbf{1} & \textbf{2} & \textbf{3} & \textbf{4} & \textbf{5} \\ \hline
\textbf{الإجابة} &  &  &  &  &  \\[8pt] \hline
\end{tabular}

\vspace{8pt}
\renewcommand{\arraystretch}{2.2}
\begin{tabular}{|>{\centering\arraybackslash}m{7.5cm}|>{\centering\arraybackslash}m{3.2cm}|>{\centering\arraybackslash}m{3.2cm}|}
\hline
\textbf{الاسم واللقب} & \textbf{الفوج} & \textbf{العلامة} \\ \hline
 & & \hfill /20 \\ \hline
\end{tabular}
\end{center}
\end{document}
